\chapter{基于特征时序缓存的视觉语言动作模型}
\label{chp:temporal_cache_vla}

\section{引言}
\label{sec:chap4_intro}

在第三章中，本文提出了一种基于扩散模型的视觉语言动作生成框架，通过跨模态注意力机制实现了视觉特征与本体状态的有效融合，并在单步动作生成任务中展现了优异的性能。然而，在真实的非结构化环境中，机器人操作任务往往具有长视距和时序依赖的特性。例如，在执行“按下按钮”等操作时，动作前后的视觉观测几乎没有差异，如果仅依赖当前单帧图像进行决策，模型极易产生“动作是否已经完成”的混淆，从而导致重复操作或任务失败。这种现象凸显了机器人操作任务本质上具有的非马尔可夫特性，即当前的决策不仅依赖于当前的观测，还需要结合历史的上下文信息。

为了解决上述长时序任务中的决策依赖问题，本章在第三章所提出的扩散动作生成模型的基础上，进一步引入了时序缓存机制，提出了一种基于特征时序缓存的视觉语言动作模型（TemporalCacheVLA）。该模型受到人类认知科学中双重记忆系统的启发，通过构建感知-语义双通道历史存储池，有效地保存和提取长视野中的低维感知细节与高维认知语义。本章将详细阐述TemporalCacheVLA框架的整体设计、视觉语言特征感知模块、时序缓存机制，以及如何将提取的时序特征接入前文所述的扩散动作生成模型中。最后，通过在LIBERO长时序任务基准上的对比实验与消融实验，系统地验证所提方法的有效性。

\section{TemporalCacheVLA框架整体设计}
\label{sec:chap4_framework}

人类在执行复杂操作任务时，通常依赖于工作记忆（Working Memory）和情景记忆（Episodic Memory）的协同作用。工作记忆负责短暂保留当前信息以供即时决策，而情景记忆则负责长期存储过去的经验细节与抽象语义。受此启发，TemporalCacheVLA框架被设计为一个“感知-存储-动作”的闭环系统，其整体架构如图~\ref{fig:temporal_cache_vla_arch}~所示。

\begin{figure}[htbp]
  \centering
  % TODO: 替换为真实的TemporalCacheVLA整体架构图路径
  \includegraphics[width=0.95\linewidth]{figures/content/temporal_cache_vla_arch}
  \caption{TemporalCacheVLA框架整体架构图}
  \label{fig:temporal_cache_vla_arch}
\end{figure}

TemporalCacheVLA框架主要由三个核心模块构成：
\begin{enumerate}
  \item \textbf{视觉语言特征感知模块}：负责接收当前的RGB图像观测与自然语言指令，利用预训练的视觉编码器和大语言模型提取细粒度的感知特征（Perceptual Tokens）和高层次的认知语义特征（Cognitive Tokens），形成当前时刻的工作记忆。
  \item \textbf{感知-语义双通道历史存储池}：作为情景记忆的载体，该模块包含感知和认知两个独立的存储通道。在每个时间步，工作记忆会通过交叉注意力机制从存储池中检索相关的历史上下文，并通过门控机制与当前特征进行自适应融合。同时，存储池会根据相似度对冗余信息进行合并压缩，以维持长时序下的紧凑表示。
  \item \textbf{基于时序缓存的扩散动作生成模块}：将融合了历史时序上下文的感知与认知特征作为条件输入，接入第三章所提出的扩散动作生成模型中。通过迭代去噪过程，生成具有时序连贯性的未来多步动作序列。
\end{enumerate}

通过上述设计，TemporalCacheVLA不仅继承了扩散模型在拟合复杂多模态动作分布方面的优势，还赋予了模型在长视距任务中理解时序依赖关系的能力。

\section{视觉语言特征感知与时序缓存机制}
\label{sec:chap4_perception_memory}

\subsection{视觉语言特征感知}
\label{subsec:chap4_perception}

在每个时间步 $t$，机器人接收当前的RGB图像观测 $I_t \in \mathbb{R}^{H \times W \times 3}$ 和自然语言指令 $L$。为了充分提取图像中的空间细节与语义信息，本文采用了双分支的视觉编码架构。具体而言，使用预训练的DINOv2\cite{oquab2023dinov2}和SigLIP\cite{zhai2023sigmoid}作为视觉主干网络，分别提取图像的局部纹理特征和全局语义特征。

提取到的原始视觉特征经过一个基于Squeeze-and-Excitation（SE）瓶颈层的感知压缩模块，被压缩为紧凑的感知特征序列 $p_t \in \mathbb{R}^{N_p \times d_p}$，其中 $N_p$ 为特征序列长度，$d_p$ 为特征维度。与此同时，原始视觉特征通过线性投影层映射到语言模型的嵌入空间中，并与分词后的语言指令拼接，输入到拥有70亿参数的大语言模型（如LLaMA\cite{touvron2023llama}）中。取大语言模型在句子结束符（EOS）位置的输出作为认知特征 $c_t \in \mathbb{R}^{1 \times d_c}$，该特征高度浓缩了当前场景的高层认知语义。

感知特征 $p_t$ 和认知特征 $c_t$ 共同构成了当前时刻的工作记忆 $M_{wk}$：
\begin{equation}
  M_{wk} = \{ p_t \in \mathbb{R}^{N_p \times d_p}, c_t \in \mathbb{R}^{1 \times d_c} \}
\end{equation}

\subsection{感知-语义双通道历史存储池}
\label{subsec:chap4_memory_pool}

为了克服单帧工作记忆缺乏时序依赖的缺陷，本文设计了感知-语义双通道历史存储池（Perceptual-Cognitive Memory Bank, PCMB）。该存储池 $M_{pcmb}$ 包含感知和认知两个通道，分别存储历史的感知细节与认知语义：
\begin{equation}
  M_{pcmb} = \{ m^x \mid x \in \{\text{per}, \text{cog}\} \}
\end{equation}
\begin{equation}
  m^x = \{ m_i^x \in \mathbb{R}^{N_x \times d_x} \}_{i=1}^L, \quad x \in \{\text{per}, \text{cog}\}
\end{equation}
其中，$L$ 为存储池的最大容量。

\textbf{记忆检索}：在每个时间步，当前的工作记忆 $M_{wk}$ 作为查询向量（Query），从存储池中检索相关的历史信息。为了保留历史特征的时序顺序，本文为存储池中的每个条目引入了正弦位置编码 $\text{TE}(\cdot)$。以感知通道为例，检索过程的键（Key）和值（Value）构建如下：
\begin{equation}
  K^x = [m_1^x + \text{TE}(t_1); \dots; m_L^x + \text{TE}(t_L)], \quad V^x = [m_1^x; \dots; m_L^x]
\end{equation}
随后，通过缩放点积注意力机制计算检索结果 $H^x$：
\begin{equation}
  H^x = \text{softmax}\left( \frac{q^x (K^x)^\top}{\sqrt{d_x}} \right) V^x, \quad q^x \in \{p_t, c_t\}, \quad x \in \{\text{per}, \text{cog}\}
\end{equation}

\textbf{门控融合}：检索到的历史特征 $H^x$ 需要与当前特征 $x$ 进行融合。为了自适应地决定历史信息的保留程度，本文设计了一个基于多层感知机（MLP）的门控机制。门控向量 $g^x$ 的计算公式为：
\begin{equation}
  g^x = \sigma(\text{MLP}(\text{concat}[x, H^x]))
\end{equation}
其中 $\sigma$ 为Sigmoid激活函数。最终的融合特征 $\tilde{x}$ 为：
\begin{equation}
  \tilde{x} = g^x \odot H^x + (1 - g^x) \odot x
\end{equation}
其中 $\odot$ 表示逐元素相乘。

\textbf{记忆压缩与更新}：融合后的特征 $\tilde{x}$ 将被更新至存储池中。当存储池的条目数量超过容量 $L$ 时，系统会计算相邻条目之间的余弦相似度，并选取相似度最高的一对相邻条目进行平均池化合并，从而在保持紧凑表示的同时剔除冗余信息：
\begin{equation}
  i^* = \arg\max_{i=1,\dots,L-1} \cos(\tilde{x}_i, \tilde{x}_{i+1})
\end{equation}
\begin{equation}
  m_{i^*}^x \leftarrow \frac{1}{2}(\tilde{x}_{i^*} + \tilde{x}_{i^*+1})
\end{equation}

\section{基于时序缓存的扩散动作生成}
\label{sec:chap4_diffusion_action}

在获取了融合时序上下文的感知特征 $\tilde{p}$ 和认知特征 $\tilde{c}$ 后，本文将其作为条件信号，接入第三章所提出的扩散动作生成模型中。第三章详细论述了扩散模型在处理多模态动作分布时的优势，以及如何通过去噪过程生成平滑的动作轨迹。本节将重点阐述时序缓存特征与扩散模型的结合方式。

假设机器人需要预测未来 $T$ 步的动作序列 $A = (a_t, a_{t+1}, \dots, a_{t+T-1})$，其中每步动作 $a_i \in \mathbb{R}^7$ 包含三维平移、三维旋转（欧拉角）以及一维的夹爪开闭状态。在扩散模型的训练阶段，真实的动作序列 $A^0$ 被逐步添加高斯噪声，生成一系列带有噪声的动作序列 $A^k$（$k \in [1, K]$ 为扩散步数）。

在去噪推理阶段，模型需要根据条件信号从纯噪声 $A^K$ 中逐步恢复出清晰的动作序列 $A^0$。本文将认知特征 $\tilde{c}$ 作为高层语义引导，将感知特征 $\tilde{p}$ 作为细粒度视觉补充。具体而言，在扩散模型的每一次去噪迭代中，噪声动作序列 $A^k$ 首先与当前去噪步数 $k$ 的正弦位置编码相加，然后与认知特征 $\tilde{c}$ 在序列维度上拼接。随后，拼接后的特征输入到基于Transformer的去噪网络中。在Transformer的交叉注意力层中，感知特征 $\tilde{p}$ 作为键（Key）和值（Value），为动作去噪提供精确的空间对齐信息。

去噪网络的优化目标是最小化预测动作与真实动作之间的均方误差（MSE）：
\begin{equation}
  \mathcal{L} = \mathbb{E}_{k, A^0, \epsilon \sim \mathcal{N}(0,I)} \left[ \| \epsilon - \epsilon_\theta(A^k, \tilde{p}, \tilde{c}, k) \|^2 \right]
\end{equation}
其中 $\epsilon_\theta$ 为去噪网络预测的噪声。通过这种方式，第三章的扩散动作生成模型被无缝集成到TemporalCacheVLA框架中，使其不仅具备了强大的多模态动作拟合能力，还获得了处理长时序依赖的全局视野。

\section{实验与分析}
\label{sec:chap4_experiments}

\subsection{实验设置}
\label{subsec:chap4_exp_setup}

为了全面评估TemporalCacheVLA在长时序和复杂操作任务中的性能，本文选取了具有高度挑战性的LIBERO\cite{liu2023libero}机器人操作基准数据集进行仿真实验。LIBERO是一个专门用于评估机器人终身学习与知识迁移能力的大规模基准，其包含了丰富的任务变体和复杂的环境交互。

LIBERO数据集共包含五个子集，分别侧重于评估机器人操作能力的不同维度：
\begin{enumerate}
  \item \textbf{LIBERO-Spatial（空间子集）}：包含10个任务，主要测试机器人在不同空间布局下的推理与操作能力，例如在不同位置抓取特定物体。
  \item \textbf{LIBERO-Object（物体子集）}：包含10个任务，侧重于评估模型对不同视觉外观、形状和纹理的物体的泛化能力。
  \item \textbf{LIBERO-Goal（目标子集）}：包含10个任务，要求机器人在相同的初始场景下，根据不同的语言指令完成不同的目标状态。
  \item \textbf{LIBERO-Long（长时序子集）}：包含10个任务，是数据集中最具挑战性的部分。该子集要求机器人执行包含多个子目标的复杂长序列操作，例如“打开抽屉，将碗放入抽屉，然后关闭抽屉”。这类任务极度依赖历史上下文信息，是验证本文时序缓存机制有效性的理想测试床。
  \item \textbf{LIBERO-90（综合子集）}：包含90个多样化的任务，全面考察模型在海量任务中的综合表现与容量极限。
\end{enumerate}

在实验中，每个任务包含50条人类专家示教轨迹。本文的模型仅使用第三人称视角的单帧RGB图像和自然语言指令作为输入，不依赖任何腕部相机视图或本体感受状态（如关节角度、末端位姿等）。模型在8张NVIDIA A100 GPU上进行训练，批大小设置为256，学习率为 $2 \times 10^{-5}$。在推理阶段，采用DDIM\cite{ho2020denoising}加速采样算法，去噪步数设置为10步。

\subsection{对比实验结果与分析}
\label{subsec:chap4_comparison}

为了验证TemporalCacheVLA的先进性，本文将其与当前领域内的主流视觉语言动作模型进行了全面的对比，包括Diffusion Policy\cite{chi2023diffusion}、Octo\cite{team2024octo}、OpenVLA\cite{kim2024openvla}、TraceVLA\cite{zheng2024tracevla}、CogACT\cite{liang2024cogact}以及$\pi_0$\cite{black2024pi}等。各方法在LIBERO五个子集上的成功率对比结果如表~\ref{tab:libero_results}~所示。

\begin{table}[htbp]
  \centering
  \caption{不同方法在LIBERO基准数据集上的成功率对比（\%）}
  \label{tab:libero_results}
  \resizebox{\textwidth}{!}{
  \begin{tabular}{lcccccc}
    \toprule
    方法 & Spatial & Object & Goal & Long & LIBERO-90 & 平均成功率 \\
    \midrule
    Diffusion Policy\cite{chi2023diffusion} & 78.3 & 92.5 & 68.3 & 50.5 & - & 72.4 \\
    Octo\cite{team2024octo} & 78.9 & 85.7 & 84.6 & 51.1 & - & 75.1 \\
    MDT\cite{reuss2024mdt} & 78.5 & 87.5 & 73.5 & 64.8 & - & 76.1 \\
    UniACT\cite{zheng2025uniact} & 77.0 & 87.0 & 77.0 & 70.0 & 73.0 & 76.8 \\
    MaIL\cite{jia2024mail} & 74.3 & 90.1 & 81.8 & 78.6 & - & 83.5 \\
    SpatialVLA\cite{qu2025spatialvla} & 88.2 & 89.9 & 78.6 & 55.5 & 46.2 & 71.7 \\
    TraceVLA\cite{zheng2024tracevla} & 84.6 & 85.2 & 75.1 & 54.1 & - & 74.8 \\
    OpenVLA\cite{kim2024openvla} & 84.7 & 88.4 & 79.2 & 53.7 & 73.5 & 75.9 \\
    CoT-VLA\cite{zhao2025cotvla} & 87.5 & 91.6 & 87.6 & 69.0 & - & 81.1 \\
    $\pi_0$-FAST\cite{pertsch2025pi0fast} & 96.4 & 96.8 & 88.6 & 60.2 & 83.1 & 85.0 \\
    TriVLA\cite{liu2025trivla} & 91.2 & 93.8 & 89.8 & 73.2 & - & 87.0 \\
    4D-VLA\cite{zhang2025_4dvla} & 88.9 & 95.2 & 90.9 & 79.1 & - & 88.6 \\
    CogACT\cite{liang2024cogact} & 97.2 & 98.0 & 90.2 & 88.8 & 92.1 & 93.2 \\
    $\pi_0$\cite{black2024pi} & 96.8 & 95.8 & 95.8 & 85.2 & - & 94.2 \\
    \midrule
    \textbf{TemporalCacheVLA (本文)} & \textbf{98.0} & \textbf{98.1} & \textbf{95.9} & \textbf{92.2} & \textbf{94.9} & \textbf{95.8} \\
    \bottomrule
  \end{tabular}
  }
\end{table}

从表~\ref{tab:libero_results}~的数据可以看出，本文提出的TemporalCacheVLA在所有子集上均取得了优异的表现，平均成功率达到了$95.8\%$，显著超越了OpenVLA（$75.9\%$）和Octo（$75.1\%$）等主流基线模型。特别值得注意的是，在最具挑战性的LIBERO-Long长时序子集上，传统缺乏时序建模的方法（如OpenVLA和Diffusion Policy）成功率仅在$50\%$左右徘徊，而TemporalCacheVLA将成功率大幅提升至$92.2\%$。这一结果强有力地证明了感知-语义双通道历史存储池在捕捉长视野时序依赖方面的关键作用。

此外，在包含90个任务的LIBERO-90综合子集上，TemporalCacheVLA同样取得了$94.9\%$的高成功率，优于CogACT的$92.1\%$和$\pi_0$-FAST的$83.1\%$。这表明本文方法不仅在处理长时序任务时具有优势，在面对大规模、多样化的任务分布时，依然能够保持强大的泛化能力和策略稳定性。

\subsection{消融实验}
\label{subsec:chap4_ablation}

为了深入分析TemporalCacheVLA中各个核心组件对整体性能的贡献，本文在长时序任务上进行了详细的消融实验。实验主要考察了时序缓存机制中记忆通道的选择、记忆检索方式以及记忆融合策略的影响。

\begin{enumerate}
  \item \textbf{记忆通道消融}：当仅使用认知语义通道时，模型能够理解任务的宏观进度，但在精细操作时容易产生偏差；当仅使用感知细节通道时，模型能够对齐空间位置，但容易在长序列中迷失当前的目标阶段。只有同时结合双通道记忆，模型才能在长时序任务中取得最佳的成功率。
  \item \textbf{位置编码消融}：在记忆检索阶段，如果去除时间步的正弦位置编码，模型将无法区分历史特征的先后顺序，导致在执行具有严格顺序要求的任务（如先开门再取物）时性能显著下降。
  \item \textbf{融合策略消融}：对比直接相加与门控融合两种策略，实验结果表明，门控机制能够更有效地过滤无关的历史噪声，自适应地提取对当前决策有用的上下文信息，从而带来更高的动作生成质量。
\end{enumerate}

\section{本章小结}
\label{sec:chap4_summary}

本章针对机器人长时序操作任务中的决策依赖问题，提出了一种基于特征时序缓存的视觉语言动作模型（TemporalCacheVLA）。该模型受人类双重记忆系统启发，设计了感知-语义双通道历史存储池，实现了对长视野中低维感知细节与高维认知语义的有效缓存与检索。同时，本章详细阐述了如何将提取的时序上下文特征作为条件信号，无缝接入第三章所提出的扩散动作生成模型中，从而赋予模型在复杂非马尔可夫环境下的全局视野与多模态动作生成能力。在LIBERO大规模基准数据集上的对比实验表明，TemporalCacheVLA在各项任务上均取得了领先的成功率，特别是在长时序任务子集上实现了显著的性能飞跃，充分验证了所提框架的有效性与先进性。
